
%----------------------------------------------------------------------------------------
\chapter{The Forecasting Engine and Material Planning Project}
\label{chap:demandforecast}
\chaptermark{Forecast Engine}
%----------------------------------------------------------------------------------------

\section{Overview of Statistical Models for Demand Forecasting}

Within a vertically integrated agri-food leader like Amadori, demand forecasting represents one of the business functions capable of generating the highest strategic 
and operational value. Operating in this sector entails managing structural complexities such as rigid biological lead times in livestock farming, strict shelf-life 
constraints on fresh poultry products, high promotional volatility in large-scale retail distribution, and the delicate necessity of balance the carcass to prevent 
costly surpluses or shortages. In this context, stochastic and statistical forecasting serves as the core predictive engine driving synchronized 
inventory management, production line balancing, and reliable procurement. Historically, statistical demand forecasting relies on decomposing time-series sales 
data into structural components: baseline level, long-term trend, recurring seasonality, and irregular noise.

Foundational methodologies range from basic moving averages to stochastic time-series models. The standard benchmark among these stochastic frameworks is the 
AutoRegressive Integrated Moving Average (ARIMA) model formalized by Box and Jenkins \cite{box1970time}, defined by parameters $p$ (autoregressive order), 
$d$ (degree of differencing to achieve stationarity), and $q$ (moving average order). Conceptually, these three parameters govern how the series processes 
historical dependency, structural trends, and unexpected disruptions: $p$ captures the memory of the process, modeling how recent sales volumes carry forward 
into the near future; $d$ addresses underlying non-stationarity and long-term trends, stabilizing the series through successive differencing; and $q$ models 
the absorption of unpredictable shocks, calibrating how transient forecasting errors from past periods dampen over time.

To accommodate predictable, cyclical demand fluctuations, this architecture expands into the Seasonal ARIMA (SARIMA) model \cite{box1970time}, 
which replicates this operational logic across a broader seasonal cycle of length $s$. The seasonal parameters mirror the baseline components: 
$P$ acts as a long-term cyclical memory, linking current consumption to the demand observed at the identical stage of previous annual cycles; 
$D$ eliminates recurring seasonal shapes by differencing observations across span $s$; and $Q$ dampens seasonal forecasting errors inherited from past cycles. 
Finally, to prevent deterministic external events---such as promotional campaigns or holiday shifts---from distorting these organic patterns, transfer function models, 
implemented as ARIMAX \cite{box1970time}, append a vector of exogenous regressors ($X$) directly to the formulation. This allows the algorithm to explicitly quantify 
and isolate external interventions, preserving the integrity of the underlying baseline series.

\section{The ``As-Is'' Demand Model: ARIMAX and TSLM Approaches}

The forecasting architecture described in this chapter reflects the operational baseline at the time of writing. 
In an enterprise of this scale, demand forecasting is subject to continuous refinement, representing an evolving process aimed at progressively optimizing 
predictive accuracy across retail dynamics. Within the current pipeline, the total demand forecast across a rolling 16-week planning horizon is generated by 
decoupling and subsequently recombining two distinct predictive streams: an ARIMAX model dedicated to estimating organic baseline demand, 
and a Time Series Linear Model (TSLM) specifically tasked with modeling promotional uplifts.

The statistical engine is executed on a weekly cadence, dynamically recalibrating model parameters across individual time series defined at the granular 
combination of Forecast Client (\textit{Cliente di Previsione}) and Article Product. 
Prior to model fitting, the historical data undergoes an automated data refactoring phase to clean the weekly sales buckets and construct the 
Reconditioned Actuals (\textit{Consuntivo Ricondizionato}) according to dedicated business logic. 
A given week is classified as a baseline period if promotional volumes account for strictly less than $33\%$ of the total historical volume recorded, 
whereas it is flagged as a promotional period if this threshold is reached or exceeded. 

The eligibility for stochastic estimation is governed by the depth of the refined time series: if a series contains at least five active weeks of Reconditioned Actuals, 
the ARIMAX algorithm is fitted; otherwise, the engine falls back onto a localized moving average. 
Due to computational constraints in processing tens of thousands of series weekly, the vector of exogenous regressors ($X$) in the current ARIMAX 
implementation is restricted exclusively to Easter calendar dynamics. 

In parallel, the TSLM model is triggered on-demand to forecast promotional spikes, leveraging promotional mechanics such as promotion type, promotional intensity, and the specific day of the week on which the event occurs.

Once computed, the forecasts are ingested into the corporate Data Warehouse, where downstream allocation routines disaggregate the figures from the aggregate Forecast Client 
down to individual Shipping Clients (\textit{Clienti di Spedizione}) and from weekly buckets into daily quantities using historical weights. 

It should be noted that the demand engine encompasses substantial computational and operational complexities that will not be further detailed, 
as an exhaustive evaluation of forecasting mathematics falls outside the scope of this thesis. 
This overview serves exclusively to contextualize the upstream demand signal, which functions as a certified, finalized input feeding directly 
into the downstream Material Planning project.


\section{The Material Planning Engine}

Downstream of the demand forecasting engine sits the Material Planning project. 
This initiative emerged from an urgent operational need to address recurring stockouts on the production lines, which historically forced planners to cut 
finished orders due to missing components. 
Because procurement lead times for ingredients, labels, and packaging can span several months from purchase order placement to delivery, 
the enterprise required an automated, data-driven system to guide purchasing and warehouse operations, replacing historical reliance on empirical manual estimation. 
The primary objective of the calculation engine is to compute the deterministic gross requirements for ingredients, packaging, and labels across both fresh and frozen product categories, 
laying the technical foundation for a future project phase that will incorporate on-hand inventory to calculate net requirements.

The calculation engine runs on a daily cadence, unifying three distinct data sources to construct a reliable, horizon-wide demand trajectory.
The short-term horizon, covering the span from the current date up to approximately ten days, ingests the operational daily production plan 
generated by the production planning system based on demand figures and refined through last minute adjustments.
The intermediate horizon, extending from the end of the production plan up to ten weeks, directly utilizes the commercial demand forecasts; 
this ten-week boundary is strictly maintained because the underlying statistical forecasting models tend to flatten out excessively over longer horizons, 
losing operational responsiveness. Beyond the eleventh week, the engine switches to the annual corporate budget; the budget represents the most approximate data source, 
as it is updated only annually and frequently contains residual volume allocations for stock-keeping units that have become obsolete throughout the fiscal year. 
For this reason, to ensure data hygiene and avoid generating fictitious material requirements, an automated upstream filtering stage strips out obsolete products.
Moreover, the engine filters out commercialized products that the enterprise merely resells without internal manufacturing, for which no Bill of Materials explosion is needed. 

A core layer of complexity stems from product taxonomy and the corresponding Bill of Materials (BOM) explosion rules across three product archetypes: 
Finished-in-Line Products (\textit{Prodotti Finiti in Linea}, PFL), Neutral Products (\textit{Prodotti Neutri}, PN), and Finished-from-Neutral Products (\textit{Prodotti Finiti da Neutro}, PFN). 
Finished-in-Line products exit the production line fully labeled and ready for the market, meaning their single BOM contains the complete set of ingredients, 
primary packaging, and labels. Neutral products represent raw, unlabeled physical items whose associated BOM includes only ingredients and primary packaging. 
Finished-from-Neutral products represent the commercial variations derived from a a Neutral Product, with a BOM containing exclusively the specific finished label(s).
Because the daily production plan schedules only PFL and PN volumes, calculating the short-term label requirements for PFN items requires computing the relative 
percentage incidence of each PFN child on its parent PN from the first ten days of the commercial forecast, reapplying these weights to the scheduled PN volume. 
Conversely, because the commercial forecast and budget plan across commercial finished codes (PFL and PFN), the engine logically maps each PFN back to its parent PN 
to explode the underlying ingredients and primary packaging.

This logic is further compounded by multi-echelon logistical routing operating within a dynamic environment characterized by frequent BOM revisions, 
platform - often referred to as deposit - modifications, and noisy master data caused by human errors.
Logistically, component requirements are allocated according to the physical stage of execution: for PFL items, the entire BOM---
encompassing ingredients, packaging, and labels - is routed directly to the production platform, where the finished product is fully manufactured, packaged, and labeled. 
Conversely, for PN items (and their downstream PFN derivatives), ingredients and packaging are exploded onto the production platform where the neutral product 
is produced and packaged, whereas the specific PFN labels are exploded onto the shipping platforms, where the final labeling operation takes place to yield the commercial PFN.

When an item can be produced or dispatched across multiple deposits, the allocation across production and shipping platforms is determined by an 
empirical routing logic designed to distribute planned volumes across the eligible sites where the Bill of Materials (BOM) is open. 
The engine first identifies all active platforms where the BOM can be exploded. If an item has an active BOM on a single platform only, 
its entire volume is assigned directly to that facility without further calculations.
When the BOM is open across multiple platforms, the engine allocates volumes based on historical incidence over the previous 12 months, 
calculating each platform's share from past volumes handled exclusively across those currently active sites. 
If operational adjustments leave an item without an active BOM, the engine falls back on past records, distributing the volume across the platforms utilized 
for that item over the past 12 months. 
For newly introduced items that lack historical data altogether, the volume is split equally across the platforms where the BOM has been opened, 
reflecting the operational rule that an item cannot be planned without an active BOM.
Both allocation processes - the one to Production Platforms and the one to Shipping Platforms - follow this identical mechanism but draw from 
domain-specific data sources: the Production Driver calculates historical shares using production batch data (\textit{lotti}), 
while the Shipping Driver derives platform percentages directly from past sales dispatch distributions. 
To ensure planning stability, these percentage shares are recalculated on a monthly cycle and frozen between iterations, remaining decoupled from daily planning executions.

Once the horizon-wide demand is fully unified and allocated across platforms, the engine executes a sequential component explosion: 
it explodes the ingredients BOM first, followed by the packaging and labeling BOM. 
The consolidated output of the engine delivers a daily, gross material requirement normalized by destination deposit, date, component code, and unit of measure, 
while maintaining full lineage traceability by preserving the detailed linkage back to the parent finished articles from which each material quantity originates. 